\SourceLeaf{52}
\[
\begin{array}{lrl}
\log(n-1)\quad .\ .\ .\ .&0.0505666&\\
\log\cos v\quad .\ .\ .\ .\ .&9.6627075&\\
\log\operatorname{tang}h.\quad .\ .\ .\ .&0.0113399&\\
\log\operatorname{cotang}(F-\delta)\quad .\ .&0.0646895&\\\cline{1-2}
\log\cos w\quad .\ .\ .\ .\ .&9.7893035,&\text{daraus }w=+52^0\ 0'\ 14''25\\
&&\text{und }u=10^0\ 36'\ 44''54
\end{array}
\]
\[
\begin{array}{lrl}
\log\operatorname{tang}h.\quad .\ .\ .\ .&0.0113399&\\
\log\cos u\quad .\ .\ .\ .\ .&9.9925074&\\\cline{1-2}
\log\operatorname{tang}G\quad .\ .\ .\ .&0.0188325,&\text{daraus }G=46^0\ 14'\ 30''77\\
&&\text{und }G-\delta=37^0\ 51'\ 47''67
\end{array}
\]
\[
\begin{array}{lrl}
\log\cos G\quad .\ .\ .\ .\ .&9.8398647&\\
\log\operatorname{tang}u.\quad .\ .\ .\ .&9.2726964&\\
\text{Komp. }\log\sin(G-\delta).&0.2119881&\\\cline{1-2}
\log\operatorname{tang}\lambda.\quad .\ .\ .\ .&9.3245492,&\text{daraus }\lambda=11^0\ 55'\ 18''31
\end{array}
\]
\[
\begin{array}{lrl}
\log\cos\lambda\quad .\ .\ .\ .\ .&9.9905300&\\
\log\operatorname{cotang}(G-\delta)\quad .\ .&0.1093281&\\\cline{1-2}
\log\operatorname{tang}\varphi\quad .\ .\ .\ .&0.0998581,&\text{daraus }\varphi=51^0\ 31'\ 47''19
\end{array}
\]

Diese Polhöhe weicht nur um wenige Bogensekunden von der wahren ab. Aus $\lambda$ ergibt sich ferner $\gamma=307^0\ 17'\ 24''91$, oder in Zeit $20^{\mathrm h}\ 29^{\mathrm m}\ 9^{\mathrm s}66$. Der Vorsprung der Uhr gegenüber der wahren Sternzeit beträgt also $=+10^{\mathrm m}\ 58^{\mathrm s}34$ und stimmt bis auf eine halbe Sekunde mit der aus anderen Beobachtungen am selben Tag gewonnenen Bestimmung überein.
\Article{7}
Unsere Untersuchung wäre gewiss unvollständig, wenn wir nicht auch zeigten, wie groß der Fehler bei der Bestimmung der unbekannten Größen werden kann, der aus den kleinen unvermeidlichen Fehlern bei der Höhenmessung entsteht. Zugleich wird daraus ersichtlich, welche Sterne man vorzugsweise wählen sollte, damit die Bestimmungen möglichst genau ausfallen. Dazu betrachten wir die Höhen $h,h'$ sowie die Größen $\varphi$ und $\lambda$ als veränderlich. Durch Differentiation von Gleichung $1$ ergibt sich dann
\expandafter\def\csname HistoricalNoteText1\endcsname{Wir setzen stets voraus, dass die Azimute vom Süden aus nach rechts über Westen, Norden usw. von $0$ bis $360^0$ gezählt werden.}
\SourceLeaf{53}
\[
\cos h\,\mathrm dh=(\sin\delta\cos\varphi-\cos\delta\sin\varphi\cos\lambda)\mathrm d\varphi-\cos\delta\cos\varphi\sin\lambda\,\mathrm d\lambda
\]
Setzt man nun das Azimut des Vertikalkreises, in dem der erste Stern beobachtet wurde, $=A$, so gilt bekanntlich\HistoricalNote{1}
\[
\begin{aligned}
\sin\delta\cos\varphi-\cos\delta\sin\varphi\cos\lambda&=-\cos h\cos A\\
\cos\delta\sin\lambda&=\cos h\sin A
\end{aligned}
\]
womit die vorangegangene Gleichung in die folgende übergeht
\[
{}[17]\qquad\mathrm dh=-\cos A\,\mathrm d\varphi-\cos\varphi\sin A\,\mathrm d\lambda
\]
Ganz ebenso ergibt sich, wenn man das Azimut des zweiten Sterns bei der zweiten Beobachtung $=A'$ setzt, durch Differentiation von Gleichung $2$
\[
{}[18]\qquad\mathrm dh'=-\cos A'\,\mathrm d\varphi-\cos\varphi\sin A'\,\mathrm d\lambda
\]
Aus der Kombination dieser Gleichungen erhält man
\[
{}[19]\qquad\mathrm d\varphi=-\frac{\sin A'}{\sin(A'-A)}\mathrm dh+\frac{\sin A\;{}'}{\sin(A'-A)}\mathrm dh'
\]
\[
{}[20]\qquad\cos\varphi\,\mathrm d\lambda=\frac{\cos A'}{\sin(A'-A)}\mathrm dh-\frac{\cos A}{\sin(A'-A)}\mathrm dh'
\]

Daraus ergibt sich unmittelbar, dass selbst die geringsten Fehler der beobachteten Höhen bei der Bestimmung sowohl der Polhöhe als auch der Zeit zu sehr beträchtlichen Fehlern führen können, wenn $\sin(A'-A)$ sehr klein ist, das heißt, wenn der zweite Vertikalkreis dem ersten oder dem ihm entgegengesetzten sehr nahe liegt. Man muss daher sorgfältig vermeiden, zwei so zueinander stehende Sterne zu kombinieren. Zugleich ist klar, dass sich die Beobachtungsfehler bei der Bestimmung der Polhöhe gewiss nur wenig verstärken können, sooft $\sin(A'-A)$ einen beträchtlichen positiven oder negativen Wert hat. Dasselbe gilt für die Zeitbestimmung, deren Genauigkeit außerdem von der Polhöhe selbst abhängt. Das liegt unabhängig von der Methode in der Natur der Sache selbst, denn in höheren Breiten lässt sich die Zeit immer mit geringerer Genauigkeit bestimmen als an Orten näher am Äquator. — Am besten wäre es stets, %
\SourceLeaf{54}die Sterne so zu wählen, dass $A'-A$ oder $A-A'$ annähernd $=90^0$ wird, dass sich also die beiden Vertikalkreise annähernd unter rechten Winkeln schneiden. In diesem Fall lässt sich leicht zeigen, dass der größte Wert von $\mathrm d\varphi$ $=\BareRoot{(\mathrm dh^2+\mathrm dh'^2)}$ ist und dass für den größten Wert der anderen Größe gilt $\mathrm d\lambda=\frac{\BareRoot{(\mathrm dh^2+\mathrm dh'^2)}}{\cos\varphi}$. Könnten wir also jede einzelne Höhe bis auf $10''$ sicher angeben, so wäre der größte Fehler bei der Bestimmung der Breite $=14''$ und der größte Fehler bei der Bestimmung des Winkels $\lambda$ wäre $\frac{14''}{\cos\varphi}$. Damit wäre in unseren Gegenden der größte Fehler bei der Zeitbestimmung $=1''5$ Zeitsekunden. Es ist kaum nötig, darauf hinzuweisen, dass bei dieser Fehlerabschätzung nur die Wirkung der Beobachtungsfehler berücksichtigt wird. Das genügt durchaus, da sich die Positionen der helleren Sterne selbst heutzutage so genau angeben lassen, dass man ihre Fehler vernachlässigen kann.

\medskip\centerline{\rule{.25\textwidth}{.3pt}}\medskip
In den Vorlesungen, die ich im kommenden Wintersemester halten werde,

\noindent\textit{öffentlich} um IX Uhr werde ich die Theorie der parabolischen Bewegung der Kometen erläutern.

\noindent\textit{privat} um X Uhr werde ich \textit{Astronomie} lehren und den Hörern zugleich Gelegenheit zu praktischen Übungen sowohl im Beobachten als auch im astronomischen Rechnen geben.
\par\medskip\centerline{\rule{.3\textwidth}{.3pt}}
